\documentclass[a4paper,10pt]{report}\input{../0.Préambule/Préambule_cours_prof}\begin{document}

\definecolor{titlecolor}{rgb}{0.98, 0.4, 0.0}
\newpage
\setcounter{chapter}{11}
\chapter{Grandeurs composées}

\section{Les grandeurs composées}

\vspace{-2mm}
\subsection{Définition}

\vspace{-2mm}
\begin{dft}
Une grandeur décrit un phénomène qui peut être mesuré ou calculé.
\end{dft}

\medskip
\noindent On distingue deux types de grandeurs :
\begin{enumerate}
\item[•] les grandeurs simples qu'on exprime à l'aide d'une unité simple :
\begin{enumerate}
\item[-] une longueur s'exprime en mètres,
\item[-] une masse s'exprime en kilogrammes,
\item[-] la capacité d'un récipient s'exprime en litres …
\end{enumerate}
\item[•] les grandeurs composées qu'on exprime avec une unité obtenue par produit ou quotient d'unités simples.
\end{enumerate}

\vspace{-2mm}
\begin{ex}
\begin{center}
\renewcommand{\arraystretch}{1.5}
\begin{tabularx}{15cm}{|*{3}{>{\centering \arraybackslash}X|}}
\hline
\rowcolor{lightgray!80} Grandeurs & Unités (symboles) & Type de grandeur \\\hline
Durée & h ; min ; s &  \\\hline
Volume & m$^3$ &  \\\hline
Intensité de courant & A (Ampère) &  \\\hline
Débit & m$^3$/s &  \\\hline
\end{tabularx}
\end{center}
\end{ex}

\vspace{-2mm}
\subsection{Les grandeurs produits}

\begin{dft}
Une grandeur produit est obtenue en multipliant deux grandeurs.
\end{dft}

\begin{ex}
L'aire d'une figure s'obtient en multipliant deux longueurs, c'est donc une grandeur produit. Elle s'exprime en : $ m ^2$
\end{ex}

\vspace{-2mm}
\begin{meth}
On détermine l'énergie électrique consommée par un appareil grâce à la formule suivante :		\quad $E = P \times t$
\begin{enumerate}
\item Quelles grandeurs désignent P et t ? Quelles sont leurs unités respectives ?

\medskip
\noindent\grandscarreaux{\linewidth}{0.8cm}

\item Déduis-en l'unité de grandeur de E.

\medskip
\noindent\grandscarreaux{\linewidth}{0.8cm}
\item Une ampoule de $40$W est restée allumée de $19$h$30$ à $23$h. Quelle énergie a-t-elle consommée ?

\medskip
\noindent\grandscarreaux{\linewidth}{2.4cm}
\end{enumerate}
\end{meth}

\vspace{-2mm}
\subsection{Les grandeurs quotients}

\vspace{-2mm}
\begin{dft}
Une grandeur quotient est obtenue en divisant deux grandeurs.
\end{dft}

\begin{ex}
La vitesse moyenne s'obtient en divisant une distance par un temps, c'est donc une grandeur quotient. 

\noindent Elle s'exprime en :	\quad $\dfrac{m}{s}$ = $m/s$ \quad ou \quad $\dfrac{km}{h}$= $km/h$
\end{ex}

\vspace{-3mm}
\begin{meth}
On mesure le débit d'un fluide grâce à la formule suivante : \quad $\text{Débit}=\dfrac{\text{volume écoulé}}{\text{temps}}$

\noindent Une canalisation fuit et perd $300$L d'eau en $20$min.
\begin{enumerate}
\item Quel est le débit de cette fuite d'eau en L/min ?

\medskip
\noindent\grandscarreaux{\linewidth}{1.6cm}
\item Convertir en m$^3$/s.

\medskip
\noindent\grandscarreaux{\linewidth}{1.6cm}
\end{enumerate}
\end{meth}

\vspace{-3mm}
\section{La notion de vitesse}

\vspace{-2mm}
\subsection{Les unités de vitesse}

\vspace{-3mm}
\noindent Les principales unités de vitesse sont :
\begin{minipage}{6cm}
\begin{center}
le mètre par seconde

\Large{m/s}
\end{center}
\end{minipage}
\begin{minipage}{6cm}
\begin{center}
le kilomètre par heure

\Large{km/h}
\end{center}
\end{minipage}

\vspace{-3mm}
\begin{ex}
En septembre 2012, Félis Baumgartner a effectué un saut d'une altitude d'environ 39 000 mètres.

\noindent Son objectif était d'être le premier homme à « dépasser le mur du son » (soit atteindre une vitesse supérieure ou égale à la vitesse du son, c'est-à-dire 340 m/s). La Fédération Aéronautique Internationale a établi qu'il avait atteint la vitesse maximale de 1 357,6 km/h au cours de sa chute libre.

\noindent A-t-il atteint son objectif ?

\medskip
\noindent\grandscarreaux{\linewidth}{1.6cm}
\end{ex}

\vspace{-2mm}
\subsection{Calculer une vitesse moyenne}

\vspace{-2mm}
\begin{dft}
La vitesse moyenne est donnée par la formule suivante : \hspace{2cm} ou bien

\medskip
\end{dft}

\begin{ex}
Le record du monde du $100$m est de $9,58$s. Quelle était la vitesse moyenne d’Usain Bolt lorsqu'il a établi ce record ?

\medskip
\noindent\grandscarreaux{\linewidth}{1.6cm}
\end{ex}

\end{document}

\begin{rmq}
par le produit en croix, on obtient :
Calcul d’une distance : d=v×t
Calcul d’un temps :  t=dv
\end{rmq}
